\documentclass{bmcart}

%%%%%%%%%%%%%%%%%%%%%%%%%%%%%%%%%%%%%%%%%%%%%%
%%                                          %%
%% CARGA DE PAQUETES DE LATEX               %%
%%                                          %%
%%%%%%%%%%%%%%%%%%%%%%%%%%%%%%%%%%%%%%%%%%%%%%

%%% Load packages
\usepackage{amsthm,amsmath}
\usepackage{graphicx}
%\RequirePackage[numbers]{natbib}
%\RequirePackage{hyperref}
\usepackage[utf8]{inputenc} %unicode support
% Español: guionado, "Tabla", "Figura", "Referencias" y "Resumen" (babel), y los
% textos que bmcart.cls deja fijos en inglés (pie de página, palabras clave, autores)
\usepackage[spanish,es-noshorthands,es-nodecimaldot,es-tabla]{babel}
\makeatletter
\def\thepage@fmt#1{Página #1 de \pageref{LastPage}}
\setattribute{keyword}{KWD}{Palabras clave}
\setattribute{corref}{text}{Correspondencia: }
\setattribute{address}{text}{Información de los autores}
\setattribute{authorinfo}{text}{La información completa de los autores está al final del artículo\endgraf}
\makeatother
%\usepackage[applemac]{inputenc} %applemac support if unicode package fails
%\usepackage[latin1]{inputenc} %UNIX support if unicode package fails


%%%%%%%%%%%%%%%%%%%%%%%%%%%%%%%%%%%%%%%%%%%%%%
%%                                          %%
%% COMIENZO DEL DOCUMENTO                   %%
%%                                          %%
%%%%%%%%%%%%%%%%%%%%%%%%%%%%%%%%%%%%%%%%%%%%%%

\begin{document}

	\begin{frontmatter}
	
		\begin{fmbox}
			\dochead{Investigación}
			
			%%%%%%%%%%%%%%%%%%%%%%%%%%%%%%%%%%%%%%%%%%%%%%
			%% INTRODUCIR TITULO PROYECTO               %%
			%%%%%%%%%%%%%%%%%%%%%%%%%%%%%%%%%%%%%%%%%%%%%%
			
			\title{Organización modular de la red de interacción de proteínas de los genes asociados a la disección aórtica (HP:0002647)}
			
			%%%%%%%%%%%%%%%%%%%%%%%%%%%%%%%%%%%%%%%%%%%%%%
			%% AUTORES. METER UNA ENTRADA AUTHOR        %%
			%% POR PERSONA                              %%
			%%%%%%%%%%%%%%%%%%%%%%%%%%%%%%%%%%%%%%%%%%%%%%
			
			\author[
			  addressref={aff1},                   % ESTA LINEA SE COPIA IGUAL PARA CADA AUTOR
			  corref={aff1},                       % ESTA LINEA SOLO DEBE TENERLA EL COORDINADOR DEL GRUPO
			  email={antoniosanchezdona9@uma.es}
			]{\inits{A.S.D.}\fnm{Antonio} \snm{Sánchez Doña}} % inits: INICIALES DE AUTOR, fnm: NOMBRE DE AUTOR, snm: APELLIDOS DE AUTOR
			\author[
			  addressref={aff1},
			  email={evalucialopezporras@uma.es}
			]{\inits{E.L.L.P.}\fnm{Eva Lucía} \snm{López Porras}}
			\author[
			  addressref={aff1},
			  email={aitananietoreus@uma.es}
			]{\inits{A.N.R.}\fnm{Aitana} \snm{Nieto Reus}}
			\author[
			  addressref={aff1},
			  email={nicosanchezalonso@uma.es}
			]{\inits{N.S.A.}\fnm{Nicolás} \snm{Sánchez Alonso}}
			
			%%%%%%%%%%%%%%%%%%%%%%%%%%%%%%%%%%%%%%%%%%%%%%
			%% AFILIACION. NO TOCAR                     %%
			%%%%%%%%%%%%%%%%%%%%%%%%%%%%%%%%%%%%%%%%%%%%%%
			
			\address[id=aff1]{%                           % unique id
			  \orgdiv{ETSI Informática},             % department, if any
			  \orgname{Universidad de Málaga},          % university, etc
			  \city{Málaga},                              % city
			  \cny{España}                                    % country
			}
		
		\end{fmbox}% comment this for two column layout
		
		\begin{abstractbox}
		
			\begin{abstract} % abstract
			
			%%%%%%%%%%%%%%%%%%%%%%%%%%%%%%%%%%%%%%%%%%%%%%%
			%% RESUMEN BREVE DE NO MAS DE 100 PALABRAS   %%
			%%%%%%%%%%%%%%%%%%%%%%%%%%%%%%%%%%%%%%%%%%%%%%%	
			
			\end{abstract}
			
			%%%%%%%%%%%%%%%%%%%%%%%%%%%%%%%%%%%%%%%%%%%%%%
			%% PALABRAS CLAVE DEL PROYECTO              %%
			%%%%%%%%%%%%%%%%%%%%%%%%%%%%%%%%%%%%%%%%%%%%%%
			
			\begin{keyword}
			\kwd{disección aórtica}
			\kwd{Human Phenotype Ontology}
			\kwd{STRING}
			\kwd{red de interacción de proteínas}
			\kwd{módulos}
			\kwd{random walk con reinicio}
			\end{keyword}
		
		
		\end{abstractbox}
	
	\end{frontmatter}
	
	%%%%%%%%%%%%%%%%%%%%%%%%%%%%%%%%%
	%% COMIENZO DEL DOCUMENTO REAL %%
	%%%%%%%%%%%%%%%%%%%%%%%%%%%%%%%%%
	
	\section{Introducción}


\subsection{La disección aórtica}

La disección aórtica es una urgencia cardiovascular grave que comienza con un desgarro de la íntima de la aorta. A través de él, la sangre penetra en la capa media y separa sus láminas, de modo que se forma una falsa luz paralela a la luz verdadera del vaso \cite{Nienaber2016}. La falsa luz puede propagarse a lo largo de la aorta, comprometer el flujo de sus ramas o romperse, lo que explica la gravedad del cuadro \cite{Nienaber2016}.

Según su extensión anatómica, la disección se describe con dos clasificaciones clásicas \cite{Isselbacher2022}. La de Stanford distingue el tipo~A, que afecta a la aorta ascendente con independencia del punto de entrada, y el tipo~B, que no la afecta. La de DeBakey diferencia el tipo~I, que se origina en la aorta ascendente y se extiende al menos hasta el arco; el tipo~II, limitado a la aorta ascendente; y el tipo~III, que se origina en la aorta descendente. Esta distinción condiciona el tratamiento: el tipo~A suele requerir cirugía urgente, mientras que el tipo~B no complicado se maneja habitualmente con tratamiento médico \cite{Isselbacher2022}.

En el estudio poblacional \textit{Oxford Vascular Study}, la incidencia de disección aórtica aguda fue de unos 6 casos por 100\,000 habitantes y año, y casi la mitad de los pacientes con disección de tipo~A (48,6\%) falleció antes de la valoración hospitalaria \cite{Howard2013}. En esa misma cohorte, aproximadamente dos tercios de los pacientes tenían hipertensión arterial conocida, con frecuencia mal controlada, lo que la sitúa como el principal factor de riesgo modificable \cite{Howard2013}. El Registro Internacional de Disección Aórtica Aguda (IRAD), que reunió más de 7300 casos en 20 años, mostró que, con la generalización de la cirugía en el tipo~A y del tratamiento endovascular en el tipo~B, la mortalidad hospitalaria se redujo de forma significativa en el tipo~A, pero no en el tipo~B \cite{Evangelista2018}.

Junto a estos factores adquiridos, una parte de los casos tiene una base hereditaria, en la que mutaciones en genes de la pared aórtica predisponen a la disección incluso en pacientes jóvenes y sin hipertensión \cite{Nienaber2016,Pinard2019}.

\subsection{Fisiopatología de la pared aórtica}

La pared de la aorta torácica está formada por tres capas, y es la capa media la que le confiere su elasticidad y su resistencia \cite{Pinard2019}. En el ser humano, esta capa se compone de más de 50 capas alternas de células de músculo liso y láminas elásticas, unidas por la denominada unidad elastina-contráctil, que conecta las láminas de elastina con el aparato contráctil de las células musculares \cite{Pinard2019}. Gracias a esta unidad, las células de músculo liso perciben las fuerzas que soporta la pared y responden a ellas \cite{Pinard2019}.

Las enfermedades aneurismáticas de la aorta comparten una degeneración proteolítica del tejido elástico y una disfunción y pérdida de células de músculo liso \cite{Pinard2019}. Se ha propuesto que las alteraciones que impiden a estas células percibir y responder a las fuerzas a través de la unidad elastina-contráctil podrían ser un motor principal de los aneurismas y las disecciones de la aorta torácica \cite{Pinard2019}. La segunda vía implicada en las formas hereditarias es la señalización del factor de crecimiento transformante beta (TGF-$\beta$): las variantes patogénicas descritas en sus genes disminuyen, o se prevé que disminuyan, la señalización celular de esta vía \cite{Pinard2019}. En el siguiente apartado se describen los genes implicados en estos procesos.

\subsection{Base genética}

Aproximadamente el 20\,\% de las personas con aneurisma o disección de aorta torácica presentan un patrón de herencia autosómica dominante en su familia, es decir, la enfermedad se debe a la mutación de un único gen \cite{Pinard2019}. Esta forma de la enfermedad se denomina enfermedad aórtica torácica hereditaria (HTAD, del inglés \textit{heritable thoracic aortic disease}) y suele transmitirse con penetrancia reducida, sobre todo en las mujeres \cite{Pinard2019}.

En algunos casos, la enfermedad aórtica forma parte de un síndrome genético \cite{Pinard2019}. El primer gen identificado fue \textit{FBN1}, que codifica la fibrilina-1 y es el causante del síndrome de Marfan. Las mutaciones en \textit{TGFBR1}, \textit{TGFBR2}, \textit{SMAD3} y \textit{TGFB2}, todos ellos de la vía de TGF-$\beta$, causan las distintas formas del síndrome de Loeys-Dietz, y las de \textit{COL3A1}, que codifica el procolágeno de tipo~III, el síndrome de Ehlers-Danlos vascular \cite{Pinard2019}. En otros pacientes la enfermedad aparece de forma aislada: el 20\,\% de los pacientes con enfermedad aórtica torácica no sindrómica tiene familiares de primer grado afectados, y solo las mutaciones en \textit{ACTA2}, que codifica la isoforma de la $\alpha$-actina propia del músculo liso, explican el 14\,\% de los casos no sindrómicos de HTAD \cite{Pinard2019}. Otros genes de este grupo codifican componentes del aparato contráctil del músculo liso, como la cadena pesada de la miosina (\textit{MYH11}) y las quinasas que regulan su contracción y relajación (\textit{MYLK} y \textit{PRKG1}) \cite{Pinard2019}. No obstante, la distinción entre formas sindrómicas y no sindrómicas es cada vez más difusa, ya que las variantes de un mismo gen pueden producir un espectro de fenotipos \cite{Pinard2019}.

Para valorar la evidencia que relaciona cada gen con la enfermedad, se aplicó el marco semicuantitativo de ClinGen a 53 genes candidatos, que se clasificaron en cinco categorías según la fuerza de su asociación; los 11 genes con evidencia definitiva o fuerte se designaron como genes de HTAD \cite{Pinard2019,Renard2018}. Según esta clasificación, los genes con evidencia definitiva son \textit{ACTA2}, \textit{COL3A1}, \textit{FBN1}, \textit{MYH11}, \textit{MYLK}, \textit{SMAD3}, \textit{TGFB2}, \textit{TGFBR1} y \textit{TGFBR2}, y los de evidencia fuerte, \textit{LOX} y \textit{PRKG1} \cite{Pinard2019,Renard2018}. Estos genes alteran tres procesos: la matriz extracelular (\textit{FBN1}, \textit{COL3A1} y \textit{LOX}), la contracción del músculo liso (\textit{ACTA2}, \textit{MYH11}, \textit{MYLK} y \textit{PRKG1}) y la señalización de TGF-$\beta$ (\textit{SMAD3}, \textit{TGFB2}, \textit{TGFBR1} y \textit{TGFBR2}) \cite{Pinard2019}. Que un mismo fenotipo pueda deberse a genes tan diversos, pero agrupados en unas pocas vías relacionadas, sugiere que estudiarlos de forma conjunta, como una red, puede ayudar a entender la enfermedad.

\subsection{El término HPO HP:0002647 (\textit{Aortic dissection})}

La Ontología de Fenotipos Humanos (HPO, del inglés \textit{Human Phenotype Ontology}) organiza y define de forma exhaustiva las características fenotípicas de las enfermedades humanas, y se ha convertido en un estándar global para la codificación computacional de datos fenotípicos detallados \cite{Gargano2024}. Al estar estructurada de forma lógica, permite realizar inferencias computacionales, y proporciona un corpus de anotaciones fenotípicas para más de 8100 enfermedades raras \cite{Gargano2024}. Por ello se utiliza en herramientas de diagnóstico de enfermedades raras, en el diagnóstico genómico y en el descubrimiento de relaciones entre genes y enfermedades \cite{Gargano2024}.

En este trabajo se partió del término HP:0002647 (\textit{Aortic dissection}) en lugar de una enfermedad concreta. Como se ha descrito, la disección aórtica aparece en varios síndromes distintos, como los de Marfan, Loeys-Dietz y Ehlers-Danlos vascular, y también en formas no sindrómicas \cite{Pinard2019}. Partir del fenotipo permitió reunir en un mismo análisis los genes asociados a la disección aórtica, con independencia de la entidad clínica en la que se describieron.


\subsection{Biología de sistemas y medicina de redes}

El fenotipo de una enfermedad rara vez depende solo del gen alterado: refleja diversos procesos patológicos que interaccionan en una red compleja \cite{Barabasi2011}. La medicina de redes se apoya en las redes de interacción de proteínas, cuyos nodos son proteínas unidas entre sí por interacciones físicas \cite{Barabasi2011}. En estas redes, las proteínas implicadas en una misma enfermedad muestran una alta tendencia a interaccionar entre sí, por lo que cabe esperar que cada enfermedad se asocie a una vecindad bien definida de la red, denominada módulo de enfermedad \cite{Barabasi2011}. Además, unos pocos nodos muy conectados (\textit{hubs}) mantienen unida toda la red \cite{Barabasi2011}. Estos enfoques se han utilizado para predecir nuevos genes de enfermedad y para identificar dianas terapéuticas y biomarcadores \cite{Barabasi2011}.

Una de las principales fuentes de interacciones entre proteínas es la base de datos STRING, que recopila, puntúa e integra información de asociación entre proteínas, tanto interacciones físicas como asociaciones funcionales \cite{Szklarczyk2025}. Sus evidencias proceden de varios canales (experimentos de laboratorio, bases de datos curadas, minería de textos de la literatura científica, coexpresión génica y contexto genómico), y cada asociación recibe una puntuación de confianza entre 0 y 1 que integra las de los distintos canales \cite{Szklarczyk2025}.

Con este enfoque, en este trabajo se construyó y se analizó la red de interacción de los genes asociados al término HP:0002647, con el fin de estudiar si estos genes se organizan en módulos funcionales y qué procesos biológicos los caracterizan. La hipótesis y los objetivos concretos se detallan en el siguiente apartado.



	\section{Hipótesis y objetivos}

\subsection*{Hipótesis}

Los genes asociados al fenotipo de disección aórtica (HP:0002647) no actúan de forma aislada, sino que se organizan en un número reducido de módulos funcionales dentro de la red de interacción de proteínas, más conectados entre sí de lo esperado por azar.

\subsection*{Objetivo general}

Caracterizar, mediante un enfoque de biología de sistemas, la red de interacción de los genes asociados al término HP:0002647 e identificar los procesos biológicos que la organizan.

\subsection*{Objetivos específicos}

\begin{enumerate}
	\item Obtener de la \textit{Human Phenotype Ontology} los genes asociados al término HP:0002647.
	\item Construir la red de interacción de estos genes a partir de STRING y describir sus propiedades topológicas, comparándolas con las de redes aleatorias del mismo tamaño.
	\item Identificar los módulos de la red mediante algoritmos de agrupamiento y evaluar su robustez frente al algoritmo utilizado y al umbral de confianza de STRING.
	\item Determinar los procesos biológicos de la lista completa de genes y de cada módulo mediante un análisis de enriquecimiento funcional.
	\item Priorizar genes candidatos asociados al fenotipo a partir de la estructura de la red.
\end{enumerate}
	\section{Materiales y métodos}

% métodos del flujo de análisis mínimo. Los números de la red salen de results/;
% las ampliaciones (A1 a A7) irán en apartados aparte cuando se suban

\subsection{Obtención de los genes asociados al fenotipo}

La lista de genes asociados al término HP:0002647 (\textit{Aortic dissection}) de la \textit{Human Phenotype Ontology} \cite{Gargano2024} se obtuvo a través de su API (\texttt{ontology.jax.org}) con el script \texttt{01\_genes.R}, que guarda los genes en \texttt{results/minimo/tablas/genes\_hpo.tsv}. La consulta se realizó el \textbf{[XX]}.
% pendiente para la ejecución final: rellenar los dos [XX] con la fecha de la consulta
% a la API de HPO (01_genes.R) y la de la descarga de genes_to_phenotype.txt

\subsection{Construcción de la red de interacción}

Los genes se buscaron en la base de datos STRING, versión 12.5 \cite{Szklarczyk2025}, mediante el paquete de R STRINGdb, versión 2.24.2 \cite{Szklarczyk2025} (script \texttt{02\_red.R}), para el ser humano (taxón 9606). Se conservaron las interacciones entre genes de la lista con una puntuación combinada igual o superior a 700, el corte de ``alta confianza'' de STRING. Este umbral se eligió tras comparar los umbrales de 400 a 900 (ampliación A1). Se eliminaron las interacciones duplicadas, se sustituyeron los identificadores de STRING por los símbolos de los genes y se registraron los genes no encontrados en STRING y los genes sin ninguna interacción con otro gen de la lista (genes aislados). Los resultados se guardaron en \texttt{results/minimo/tablas/red.tsv} y \texttt{genes\_string.tsv}.

\subsection{Análisis de la topología de la red}

La red se analizó como un grafo no dirigido con el paquete de R igraph, versión 2.3.4 \cite{Csardi2006,Antonov2023,Csardi2026} (script \texttt{03\_topologia.R}). Se incluyeron todos los genes encontrados en STRING, con los genes aislados como nodos sueltos, de modo que cuentan en el número de nodos y en la densidad. Se calcularon el número de nodos y de aristas, el grado medio, la densidad, el número de componentes conexas y el tamaño de la mayor. La densidad es la proporción de parejas de genes conectadas entre las posibles; se calculó con todos los nodos (75 aristas de 820 parejas posibles) y, para comparar, sin los genes aislados (75 de 465 parejas). Como la red no es conexa, la cercanía habitual no es comparable entre componentes. Por eso se calculó la cercanía dentro de la componente mayor y la centralidad armónica para toda la red. Los resultados se guardaron en \texttt{results/minimo/tablas/topologia\_resumen.tsv} y \texttt{topologia\_genes.tsv}.

\subsection{Detección de módulos}

Los módulos de la red se identificaron con el algoritmo de Louvain \cite{Blondel2008}, implementado en igraph (script \texttt{04\_clustering.R}). Como el algoritmo procesa los nodos en orden aleatorio, se fijó la semilla del generador de números aleatorios en 123 para que el resultado sea reproducible. El análisis se hizo sin los genes aislados, porque cada uno formaría un módulo de un solo gen. Al incluirlos se obtienen los mismos módulos y la misma modularidad. La asignación de genes a módulos se guardó en \texttt{results/minimo/tablas/clusters.tsv}.

\subsection{Enriquecimiento funcional}

El enriquecimiento funcional se realizó con la función \texttt{enrichGO} del paquete de R clusterProfiler, versión 4.20.3 \cite{Wu2021}, con las anotaciones de \texttt{org.Hs.eg.db} y los términos de procesos biológicos de Gene Ontology \cite{Ashburner2000} (script \texttt{05\_enriquecimiento.R}). El análisis se basó en un test de sobrerrepresentación (hipergeométrico) y los valores de $p$ se ajustaron por el método de Benjamini-Hochberg; se consideraron significativos los términos con $p$ ajustado menor que 0,05 y valor $q$ menor que 0,2. Para la lista completa de genes se usó como fondo el conjunto de genes con alguna anotación en la HPO (archivo \texttt{genes\_to\_phenotype.txt}, descargado el \textbf{[XX]}), porque el fondo por defecto, todos los genes con anotación GO, infla el número de términos. Solo se enriquecieron los módulos con al menos 5 genes, ya que con 2 o 3 genes aparecen términos significativos sin interpretación biológica, y para ellos el fondo fueron los genes de la red. Las tablas se guardaron en \texttt{results/minimo/tablas/} (archivos \texttt{GO\_*.csv}).

\subsection{Ampliación A1: barrido de umbrales de STRING}

Para estudiar cómo cambia la red según la confianza que se exige a las interacciones, se construyó la red con seis umbrales de puntuación combinada de STRING (400, 500, 600, 700, 800 y 900), cada uno con el mismo script \texttt{02\_red.R} del análisis mínimo, y se compararon con el script \texttt{A1\_barrido\_umbrales.R}. De cada umbral se calcularon el número de genes con y sin interacción, el número de aristas, la densidad (con los genes aislados como nodos sueltos, como en el análisis de la topología), el tamaño de la componente mayor, el número de módulos de Louvain, el número de módulos de al menos 5 genes y la modularidad. Los módulos se detectaron sin los genes aislados y con la misma semilla (123) que en el análisis mínimo. Además, para cada gen se anotó si tiene alguna interacción con cada umbral y el umbral más alto con el que todavía la tiene. Los resultados se guardaron en \texttt{results/ampliaciones/A1\_umbrales/} (archivos \texttt{barrido\_umbrales.tsv}, \texttt{genes\_por\_umbral.tsv} y \texttt{barrido\_umbrales.pdf}).

\subsection{Ampliación A2: red real frente a redes aleatorias}

Para saber si la organización en módulos de la red del análisis mínimo es mayor de lo que cabría esperar por azar, se comparó la red real (no dirigida y sin los genes aislados) con redes aleatorias de dos modelos (script \texttt{A2\_redes\_aleatorias.R}). El primero fue el modelo de Erd\H{o}s--R\'enyi (función \texttt{sample\_gnm} de igraph), que coloca al azar el mismo número de interacciones entre el mismo número de genes. El segundo conservó el grado: se reordenaron las interacciones de la red real (función \texttt{rewire} con \texttt{keeping\_degseq}) de manera que cada gen mantuvo su número de interacciones y solo cambió con qué genes interacciona. De cada modelo se generaron 1000 redes.

En la red real y en cada red aleatoria se calcularon cuatro propiedades: la modularidad de Louvain, la transitividad global, la distancia media entre los pares de genes conectados y la asortatividad por grado. Para cada propiedad y modelo se calculó el valor $z$ (Ecuación~\ref{eq:z}) y un valor $p$ empírico (Ecuación~\ref{eq:p}):
\begin{equation}
	z = \frac{x_{\mathrm{real}} - \mu_{\mathrm{azar}}}{\sigma_{\mathrm{azar}}}
	\label{eq:z}
\end{equation}
\begin{equation}
	p = \frac{1 + k}{1 + n}
	\label{eq:p}
\end{equation}
donde $x_{\mathrm{real}}$ es el valor de la propiedad en la red real, $\mu_{\mathrm{azar}}$ y $\sigma_{\mathrm{azar}}$ son la media y la desviación típica de esa propiedad en las redes aleatorias, $n$ es el número de redes aleatorias (1000) y $k$ es el número de ellas con un valor igual o más extremo que el real, en la dirección en que este se aparta de la media. Con 1000 redes, el menor valor posible de $p$ es $1/1001 \approx 0{,}001$. Todos los procesos con azar usaron la semilla 123. Los resultados se guardaron en \texttt{results/ampliaciones/A2\_redes\_aleatorias/} (archivos \texttt{A2\_nulos.tsv}, \texttt{A2\_resumen.tsv} y \texttt{A2\_redes\_aleatorias.pdf}).


\subsection{Ampliación A3: ataque a los hubs frente a fallo al azar}

Para estudiar la robustez de la red del análisis mínimo (no dirigida y sin los genes aislados), se clasificaron los genes según su grado y su intermediación y se midió cuánto se rompía la red al quitarlos (script \texttt{A3\_ataque\_hubs.R}). El grado es el número de interacciones de un gen, y la intermediación (\textit{betweenness}) es el número de caminos más cortos entre otros dos genes que pasan por él, ambos se calcularon con igraph. Se consideraron hubs los genes con el 20 \% de grado más alto y cuellos de botella los genes con el 20 \% de intermediación más alta, con el mismo corte para los dos, como en el trabajo de Yu \textit{et al.} \cite{Yu2007}. De cada gen se anotó también su módulo de Louvain y el número de módulos distintos entre sus vecinos.

Después se quitaron los genes de uno en uno y, tras cada paso, se midió el tamaño de la componente mayor de la red \cite{Albert2000}. Se siguieron tres órdenes: de mayor a menor grado (ataque por grado), de mayor a menor intermediación (ataque por intermediación) y 1000 órdenes aleatorios (fallo al azar, con la semilla 123). Los empates se resolvieron por el nombre del gen. Se compararon los tres órdenes al quitar el 20 \% de los genes (6 de 31) y según el número de genes que hay que quitar para que la componente mayor quede en la mitad o menos de su tamaño inicial. El valor $p$ empírico se calculó como en la ampliación A2 (Ecuación~\ref{eq:p}), tomando como $k$ el número de órdenes aleatorios que dejan una componente mayor igual o más pequeña que el ataque (o que la parten con tan pocos genes como él, o menos). Por último, se calculó la conectividad de la componente mayor (funciones \texttt{vertex\_connectivity} y \texttt{edge\_connectivity} de igraph), que es el número mínimo de genes o de interacciones que hay que quitar para partirla. Los resultados se guardaron en \texttt{results/ampliaciones/A3\_ataque\_hubs/} (archivos \texttt{A3\_hubs\_cuellos.tsv}, \texttt{A3\_curvas.tsv}, \texttt{A3\_resumen.tsv} y \texttt{A3\_ataque\_hubs.pdf}).

\subsection{Ampliación A4: otros algoritmos de clustering}

Para comprobar si los módulos dependen del algoritmo elegido, se aplicaron a la red del análisis mínimo (no dirigida y sin los genes aislados) siete algoritmos de detección de comunidades de igraph: Louvain \cite{Blondel2008}, Leiden \cite{Traag2019} (con la modularidad como función objetivo), \textit{fast greedy} \cite{Clauset2004}, \textit{walktrap} \cite{Pons2006}, \textit{infomap} \cite{Rosvall2008}, \textit{label propagation} \cite{Raghavan2007} y \textit{edge betweenness} \cite{Girvan2002} (script \texttt{A4\_clusterings.R}). De cada algoritmo se calcularon el número de módulos, su modularidad y el tamaño de cada módulo. Las particiones se compararon con la información mutua normalizada (NMI) \cite{Danon2005}, que vale 1 si dos algoritmos agrupan los genes de la misma manera y se acerca a 0 cuanto menos se parecen; se calculó la NMI de cada algoritmo frente a Louvain y la de cada pareja de algoritmos. También se aplicó Louvain con tres valores del parámetro \texttt{resolution} (0,5, 1 y 1,5), donde 1 es el valor por defecto usado en el análisis mínimo, para ver si el número de módulos depende de este parámetro. Por último, los siete algoritmos se aplicaron a la red de cada umbral de la ampliación A1, para ver si los módulos dependen del umbral además del algoritmo. Todos los procesos con azar usaron la semilla 123. Los resultados se guardaron en \texttt{results/ampliaciones/A4\_clusterings/} (tablas \texttt{A4\_*.tsv} y figura \texttt{A4\_clusterings.pdf}).

\subsection{Reproducibilidad}

Todo el análisis se programó en R (versión 4.6.2) y bash. El script \texttt{setup.sh} instala los paquetes dentro del repositorio, sin permisos de administrador, y \texttt{launch.sh} ejecuta el flujo completo y guarda todo lo que genera en \texttt{results/}. Los procesos con azar usaron la semilla 123.
	\section{Resultados}

\subsection{Propiedades de la red}
La red de interacción de STRING (versión 12.5, umbral 700) de los genes asociados a HP:0002647 se construyó con todos los genes encontrados en STRING, de modo que los genes sin ninguna interacción quedaron como nodos aislados. La Tabla~\ref{tab:red} resume sus propiedades (archivo \texttt{results/minimo/tablas/topologia\_resumen.tsv}). Diez de los 41 genes quedaron aislados: \textit{SLC2A10}, \textit{IPO8}, \textit{FOXE3}, \textit{PRKG1}, \textit{MAT2A}, \textit{THSD4}, \textit{ANGPTL6}, \textit{THSD1}, \textit{ROBO4} y \textit{GATA5}. Los genes con más interacciones fueron \textit{TGFBR1} y \textit{TGFBR2} (11 interacciones cada uno) y \textit{FBN1} (10) (archivo \texttt{topologia\_genes.tsv}).

\begin{table}[h!]
	\renewcommand{\arraystretch}{1.25}
	\caption{Propiedades de la red de interacción de los genes asociados a HP:0002647 (STRING 12.5, umbral 700). Fuente: \texttt{results/minimo/tablas/topologia\_resumen.tsv}.}
	\label{tab:red}
	\begin{tabular}{lrp{6.3cm}}
		\hline
		Propiedad & Valor & Interpretación \tabularnewline
		\hline
		Nodos & 41 & Todos los genes de HPO se encontraron en STRING \tabularnewline
		Nodos con interacción & 31 & Genes con al menos una interacción con otro gen de la lista \tabularnewline
		Nodos aislados & 10 & Genes sin ninguna interacción con otro gen de la lista \tabularnewline
		Aristas & 75 & Interacciones entre pares de genes \tabularnewline
		Grado medio & 3,659 & Cada gen interacciona, de media, con unos 3,7 genes de la lista \tabularnewline
		Densidad & 0,0915 & Proporción de las 820 parejas posibles que están conectadas \tabularnewline
		Componentes conexas & 13 & La red no es conexa: tiene 13 bloques, entre ellos los 10 genes aislados \tabularnewline
		Nodos de la componente mayor & 26 & 26 de los 31 genes conectados forman un único bloque \tabularnewline
		\hline
	\end{tabular}
\end{table}
	\section{Discusión}

\subsection{Organización modular de la red}

Los genes asociados al término HP:0002647 no se distribuyeron al azar en la red de interacción, sino que se agruparon en seis módulos, dos de ellos de mayor tamaño. Esta organización concuerda con el concepto de módulo de enfermedad, según el cual las proteínas implicadas en un mismo fenotipo tienden a interaccionar entre sí y a ocupar una vecindad concreta de la red \cite{Barabasi2011}. Además, los módulos reprodujeron los tres procesos que la literatura relaciona con la enfermedad aórtica torácica hereditaria: la matriz extracelular, la señalización de TGF-$\beta$ y la contracción del músculo liso \cite{Pinard2019,Lindsay2011}. Este resultado es, por tanto, el esperado y sirve sobre todo para confirmar que la red recoge la biología conocida de la enfermedad. Su interés reside más bien en cómo se relacionan estos procesos entre sí y en los módulos que no encajan en este esquema.

Los dos módulos principales, de diez genes cada uno, correspondieron a la matriz extracelular y las fibras elásticas (\textit{BGN}, \textit{COL3A1}, \textit{EFEMP2}, \textit{ELN}, \textit{EMILIN1}, \textit{FBN1}, \textit{LOX}, \textit{MFAP5}, \textit{P4HA2} y \textit{PLOD1}) y a la señalización de TGF-$\beta$/SMAD (\textit{ENG}, \textit{SMAD2}, \textit{SMAD3}, \textit{SMAD4}, \textit{SMAD6}, \textit{TGFB2}, \textit{TGFB3}, \textit{TGFBR1}, \textit{TGFBR2} y \textit{TGFBR3}). Aunque el algoritmo de Louvain los separó, ambos procesos están estrechamente ligados: los modelos animales de las formas hereditarias mostraron la implicación inesperada de la vía de TGF-$\beta$ en enfermedades causadas por defectos de proteínas estructurales de la matriz, como la fibrilina-1 en el síndrome de Marfan \cite{Lindsay2011}. En consonancia con ello, los tres genes con más interacciones de la red, \textit{TGFBR1}, \textit{TGFBR2} y \textit{FBN1}, pertenecieron a estos dos módulos, lo que sitúa a la matriz extracelular y a la señalización de TGF-$\beta$ en el núcleo de la red.

El módulo de músculo liso agrupó a \textit{ACTA2}, \textit{MYH11} y \textit{MYLK}, que codifican componentes del aparato contráctil y de su regulación. Estas proteínas forman parte de la unidad elastina-contráctil, la estructura que une las fibras elásticas con las células de músculo liso de la capa media y cuya alteración es común a muchos de los genes que causan aneurismas y disecciones de la aorta torácica \cite{Karimi2016}. Llama la atención que \textit{PRKG1}, que codifica una quinasa de la misma vía y cuenta con evidencia fuerte de asociación con la enfermedad \cite{Pinard2019,Renard2018} quedara sin interacciones con el resto de genes con el umbral utilizado. Este caso muestra que la posición de un gen en la red no refleja necesariamente su relevancia clínica.

El módulo de Notch (\textit{NOTCH1}, \textit{HEY2} y \textit{NKX2-5}) apunta a un mecanismo distinto, relacionado con el desarrollo cardiovascular más que con el mantenimiento de la pared aórtica. Las mutaciones en \textit{NOTCH1} causan un espectro de anomalías del desarrollo de la válvula aórtica, y los represores transcripcionales de la familia de \textit{HEY2}, activados por la señalización de Notch1, intervienen en ese mismo proceso \cite{Garg2005}.

\subsection{Módulos inesperados}

El módulo formado por \textit{HLA-B}, \textit{HLA-DRB1} y \textit{PTPN22} fue el resultado más inesperado, ya que reúne genes del sistema inmunitario sin relación con la estructura de la pared aórtica. Según las anotaciones de HPO, los tres genes se asociaron al término a través de una misma enfermedad, la arteritis de células gigantes (ORPHA:397), una vasculitis de grandes vasos. En esta enfermedad, la región HLA concentra la mayor parte de la susceptibilidad genética, con contribuciones de \textit{HLA-DRB1} y de \textit{HLA-B} \cite{Carmona2015}, y la variante R620W de \textit{PTPN22} constituye otro factor de riesgo \cite{Serrano2013}. Que este módulo no se conectara con los de matriz extracelular y TGF-$\beta$ es coherente con una vía de daño de la pared aórtica distinta, de origen inflamatorio y no estructural.

Del mismo modo, \textit{ABCC6} y \textit{ENPP1}, ambos asociados al término a través de la calcificación arterial generalizada de la infancia (ORPHA:51608), formaron un módulo aislado del resto. Las mutaciones en cualquiera de los dos genes pueden causar la calcificación arterial generalizada de la infancia y el pseudoxantoma elástico, dos extremos de un mismo espectro de calcificación ectópica que afecta a la lámina elástica interna de las arterias \cite{Nitschke2012}. Este módulo representaría, por tanto, una tercera vía de daño de las fibras elásticas, a través de su calcificación.

\subsection{Enriquecimiento funcional y elección del fondo}

El análisis de enriquecimiento de la lista completa, con el fondo de los 5277 genes con anotación en HPO, identificó 410 procesos biológicos de GO, encabezados por \textit{extracellular matrix assembly}, y el del módulo de TGF-$\beta$, 81 procesos, entre ellos \textit{SMAD protein signal transduction}. Ambos resultados confirmaron la interpretación biológica de los módulos.

En cambio, el módulo de matriz extracelular no mostró ningún término enriquecido. Este resultado no indica que el módulo carezca de una función común, sino que se debe al fondo utilizado: en el análisis de cada módulo, el fondo fueron los 31 genes de la red, que ya son en buena parte proteínas de la matriz extracelular. Como el test de enriquecimiento mide si una función está más representada en la lista que en el fondo, un módulo cuya función coincide con la del fondo no puede destacar. Esto ilustra que la elección del fondo condiciona el resultado del enriquecimiento, una fuente de sesgo descrita en la literatura que a menudo se pasa por alto \cite{Timmons2015}. Por el mismo motivo, para la lista completa se utilizó como fondo el conjunto de genes con alguna anotación en HPO y no todos los genes con anotación en GO: como la lista procede de HPO, compararla con todo el genoma habría inflado el número de términos significativos. Un número menor de términos con el fondo adecuado es, por tanto, un resultado más fiable.

Por último, solo se enriquecieron los módulos de cinco genes o más, ya que en listas tan pequeñas unos pocos genes bastan para que un término resulte significativo. Por ello, los módulos pequeños (músculo liso, Notch, HLA/\textit{PTPN22} y \textit{ABCC6}/\textit{ENPP1}) se interpretaron únicamente a partir de la literatura.
	\input{tex_files/conclusiones.tex}
	
	
	%%%%%%%%%%%%%%%%%%%%%%%%%%%%%%%%%%%%%%%%%%%%%%
	%% OTRA INFORMACIÓN                         %%
	%%%%%%%%%%%%%%%%%%%%%%%%%%%%%%%%%%%%%%%%%%%%%%
	
	\begin{backmatter}
	
		\section*{Abreviaciones}%% if any
			Indicar lista de abreviaciones mostrando cada acrónimo a que corresponde
		
		\section*{Disponibilidad de datos y materiales}%% if any
			El código, con las instrucciones para reproducir todos los resultados, está disponible en \texttt{https://github.com/antoniodoa9/project\_template}.
		
		\section*{Contribución de los autores}
			Usando las iniciales que habéis definido al comienzo del documento, debeis indicar la contribución al proyecto en el estilo:
			J.E : Encargado del análisis de coexpresión con R, escritura de resultados; J.R.S : modelado de red con python y automatizado del código, escritura de métodos; ...
			OJO: que sea realista con los registros que hay en vuestros repositorios de github. 
		
		
		%%%%%%%%%%%%%%%%%%%%%%%%%%%%%%%%%%%%%%%%%%%%%%%%%%%%%%%%%%%%%%%%%%%%%%%%%%%%%%%%%%%%%%%%
		%% BIBLIOGRAFIA: no teneis que tocar nada, solo sustituir el archivo bibliography.bib %%
		%% por el que hayais generado vosotros                                                %%
		%%%%%%%%%%%%%%%%%%%%%%%%%%%%%%%%%%%%%%%%%%%%%%%%%%%%%%%%%%%%%%%%%%%%%%%%%%%%%%%%%%%%%%%%
		
		\bibliographystyle{bmc-mathphys} % Style BST file (bmc-mathphys, vancouver, spbasic).
		\bibliography{bibliography}      % Bibliography file (usually '*.bib' )
	
	\end{backmatter}
	
	\clearpage
\section*{Material suplementario}
\setcounter{figure}{0}
\renewcommand{\thefigure}{S\arabic{figure}}
\setcounter{table}{0}
\renewcommand{\thetable}{S\arabic{table}}

	
\end{document}
